% Part: turing-machines
% Chapter: machines-computations
% Section: variants

\documentclass[../../../include/open-logic-section]{subfiles}

\begin{document}

\olfileid[tr]{tur}{mac}{var}
\olsection{Turing Makinesi Çeşitlemeleri}

Aslında Turing makinelerini tanımlamanın birçok olası yolu vardır; bizimki
bunlardan yalnızca biridir. Tanımımız bazı bakımlardan ötekilerden daha
serbesttir. Keyfî sonlu alfabelere izin veriyoruz; daha kısıtlı bir tanım
yalnızca $\TMstroke$ ve~$\TMblank$ olmak üzere iki şerit simgesine izin
verebilir. Makinenin şeride bir simge yazıp aynı anda hareket etmesine izin
veriyoruz; başka tanımlar ya yazmaya ya da hareket etmeye izin verir.
Makinenin şerit kafasını oynatmadan yazmasına izin veriyoruz; başka tanımlar
$\TMstay$ ``yönergesini'' içermez. Tanımımız başka bakımlardan daha
kısıtlayıcıdır. Şeridin yalnızca bir yönde sonsuz olduğunu varsaydık; başka
tanımlar şeridin hem sola hem sağa sonsuz olmasına izin verir. Hatta herhangi
sayıda ayrı şeride ya da sonsuz bir kareler ızgarasına bile izin verilebilir.
Turing makinesinin yönerge kümesini bir geçiş fonksiyonuyla gösteriyoruz;
başka tanımlar, makinenin belirli bir durumda birden fazla olası yönergeye
sahip olduğu bir geçiş bağıntısı kullanır.

Tanımın bu son gevşetilişi özellikle ilginçtir. Bizim tanımımızda makine
$q$ durumunda $\sigma$ simgesini okurken $\delta(q,\sigma)$ yeni simgeyi,
durumu ve şerit kafası konumunu belirler. Ancak yönerge kümesinin mevcut
durum-simge çiftleri $\tuple{q,\sigma}$ ile yeni durum-simge-yön üçlüleri
$\tuple{q',\sigma',D}$ arasında bir bağıntı olmasına izin verirsek Turing
makinesinin eylemi biricik biçimde belirlenmeyebilir: Yönerge bağıntısı hem
$\tuple{q,\sigma,q',\sigma',D}$ hem de
$\tuple{q,\sigma,q'',\sigma'',D'}$ içerebilir. Bu durumda
\emph{belirlenimsiz} bir Turing makinemiz vardır. Bu makineler hesaplama
karmaşıklığı kuramında önemli bir rol oynar.

Bir Turing makinesinin ne zaman durduğuna ilişkin de farklı uzlaşımlar vardır:
Biz geçiş fonksiyonu tanımsız olduğunda durduğunu söyleriz; başka tanımlar
makinenin özel olarak belirlenmiş bir durma durumunda olmasını gerektirir.
Belirli bir durma durumu istemenin Turing makinelerinin hesaplayabileceklerini
etkileyen bir kısıtlama olmadığını \olref[hal]{sec} içinde açıkladık. Turing
makinelerimizin şeritleri yalnızca bir yönde sonsuz olduğundan bir Turing
makinesinin bir yönergeyi gerektiği gibi yerine getiremeyeceği durumlar vardır:
En soldaki kareyi okuyup sola gitmesi gerekiyorsa. Tanımımıza göre ``şeritten
düşmek'' yerine yerinde kalır; ama böyle olduğunda duracak biçimde de
tanımlayabilirdik. Bu tanım da eşdeğerdir: $0$. kareden sola gitmeye
çalıştığında duran bir Turing makinesinin davranışını, her
$\delta(q,\TMendtape)=\tuple{q',\sigma,\TMleft}$ geçişini silerek
benzetebiliriz; böylece makine $\TMendtape$ üzerinde sola gitmeye çalışmak
yerine durur.\footnote{Bu çözüm \emph{tam olarak} işlemez; çünkü
$\TMendtape$ simgesini $0$. kare dışındaki karelere yazmamızı ve oralarda
okumamızı hiçbir şey engellemez (bkz.~\olref[una]{ex:mover}). Böyle bir amaç
için kullanılacak ikinci bir~$\TMendtape'$ simgesi ekleyerek bunu aşabiliriz.}

Sayıları (dolayısıyla bir Turing makinesinin hesapladığı girdi-çıktı
fonksiyonunu) göstermenin de farklı yolları vardır: Biz tekli gösterimi
kullanıyoruz, fakat ikili gösterim de kullanılabilir. Bunun için $\TMblank$
ve~$\TMendtape$ simgelerine ek olarak iki simge gerekir.

İşte ilginç bir olgu: Bu çeşitlemelerin hiçbiri hangi fonksiyonların Turing
hesaplanabilir olduğunu değiştirmez. \emph{Bir fonksiyon bir tanıma göre
Turing hesaplanabilir ise bütün bu tanımlara göre Turing hesaplanabilirdir.}

Bunu doğrulamanın ayrıntılarına girmeyeceğiz. Yalnızca bir örnek verelim:
Her iki yönde sonsuz bir şeride ya da birden fazla şeride izin vermek ek
hesaplama gücü kazandırmaz. Kabaca nedeni, tek yönlü sonsuz tek şeritli bir
Turing makinesinin çoklu ya da iki yönlü sonsuz şeritleri benzetebilmesidir.
Örneğin ek durumlar ve yönergeler kullanarak çok şeritli ya da iki yönlü
sonsuz şeritli bir makinenin programını tek yönlü sonsuz tek şeritli bir
makineye ``çevirebiliriz''. Çevrilmiş makine çift numaralı kareleri $1$.
şeridin kareleri (ya da iki yönlü sonsuz bir şeridin ``pozitif'' kareleri),
tek numaralı kareleri ise $2$. şeridin kareleri (ya da ``negatif'' kareler)
için kullanabilir.

\end{document}
